\section{Operator norms in orthonormal coordinates}

\begin{theorem}[Operator norm in orthonormal coordinates]
    \label{thm:operator_norm_orthonormal_coordinates}
    \lean{ContinuousLinearMap.norm_toMatrix_orthonormal}
    \leanok
    Let $E$ and $F$ be inner-product spaces over $\mathbb K=\R$ or $\C$,
    with finite orthonormal bases $(e_j)_{j\in J}$ and $(f_i)_{i\in I}$.
    For a continuous linear map $T:E\to F$, let
    $M_{ij}=\langle f_i,Te_j\rangle$ be its matrix in these bases.
    The Euclidean operator norm of $M$ equals the operator norm of $T$:
    $\|M\|_{2\to2}=\|T\|$.
    The two basis index sets may have different cardinalities, and either
    may be empty.
\end{theorem}

\begin{proof}\leanok
    Let $U:E\to\mathbb K^J$ and $V:F\to\mathbb K^I$ be the coordinate
    isometries. The linear map represented by $M$ is $VTU^{-1}$.
    Composition on either side with a surjective linear isometry preserves
    the operator norm, so $\|M\|_{2\to2}=\|VTU^{-1}\|=\|T\|$.
\end{proof}

\begin{theorem}[Expectations under a change of coordinates]
    \label{thm:expectation_coordinate_change}
    \lean{Matrix.inner_toEuclideanLin_reindex}
    \leanok
    Let $X$ and $Y$ be finite sets, $e:X\to Y$ a bijection,
    $M\in\C^{X\times X}$ and $\psi\in\C^X$.
    Define $U_e:\C^X\to\C^Y$ by $(U_e\psi)(y)=\psi(e^{-1}y)$, and define
    $M_e\in\C^{Y\times Y}$ by
    $(M_e)_{y,y'}=M_{e^{-1}y,e^{-1}y'}$.
    Then $\langle U_e\psi,M_eU_e\psi\rangle=\langle\psi,M\psi\rangle$.
    The index sets may be empty.
\end{theorem}

\begin{proof}\leanok
    Relabelling preserves the inner product, and matrix multiplication
    gives $M_eU_e\psi=U_e(M\psi)$. Applying the inner-product identity to
    $\psi$ and $M\psi$ proves the assertion.
\end{proof}

\begin{theorem}[Operator norm after adjoining an identity factor]
    \label{thm:operator_norm_identity_factor}
    \lean{Matrix.l2_opNorm_one_kronecker}
    \leanok
    Let $X$ and $Y$ be finite sets with $X\ne\varnothing$, and let
    $K\in\C^{Y\times Y}$.
    For the Euclidean operator norm,
    $\lVert I_X\otimes K\rVert=\lVert K\rVert$.
    The set $Y$ may be empty.
\end{theorem}

\begin{proof}\leanok
    The map $\Phi:K\mapsto I_X\otimes K$ is a complex star-algebra
    homomorphism. Choose $x_0\in X$.
    Since $\Phi(K)_{(x_0,i),(x_0,j)}=K_{ij}$, the map is injective.
    An injective star-algebra homomorphism between complex C-star algebras
    is isometric, giving the claimed equality.
\end{proof}
